% 01-intro — the 1990→2026 mission statement
\section{Introduction}
\label{sec:intro}

In 1990, Hornbostel, Brodsky and Pauli diagonalized the light-cone
Hamiltonian of SU($N$) gauge theory in one space and one time
dimension~\cite{Hornbostel:1988fb}.  They computed the hadronic spectrum and
wavefunctions across couplings, colors, and baryon number at resolutions
$2K = 16$--$24$, in minutes on the mainframes of the day.  Convergence was
controlled by a single extrapolation formula.  The paper became the standard
reference for discretized light-cone quantization (DLCQ), and its figures
are still the picture most physicists carry of QCD in two dimensions at
finite $N$.

The method has a long lineage.  DLCQ was introduced by Pauli and
Brodsky~\cite{Pauli:1985pv,Pauli:1985ps} and first applied to the Schwinger
model~\cite{Eller:1986nt}; reviews are
Refs.~\cite{Brodsky:1992mp,Brodsky:1997de}.  The same framework later
reached finite temperature in QED$_{1+1}$, including thermodynamics and the
analogue of the $R$ ratio~\cite{Strauss:2008zx,Strauss:2009zza}.  Light-cone
wavefunctions feed a much wider class of observables than spectra.

Nor did the program stop in two dimensions.  Positronium and quarkonium
served as 3+1 DLCQ test cases early~\cite{Krautgartner:1991xz}, front-form
QED resolved the positronium spin
multiplets~\cite{Trittmann:1997up,Trittmann:1997tt}, and parallel
light-front solvers later delivered precision true-muonium
spectra~\cite{Lamm:2013oga,Lamm:2016ong}.  Higher-dimensional QCD advanced
through transverse-lattice methods~\cite{Burkardt:2001jg}, supersymmetric
DLCQ and its relatives~\cite{Hiller:2016itl}, and basis light-front
quantization~\cite{Vary:2009gt}, which now produces quarkonium spectra in
full 3+1 dimensions~\cite{Li:2015zda}.  Two-dimensional adjoint QCD remains
under active light-front study with basis-function
methods~\cite{Trittmann:2023dar}.  Ref.~\cite{Bakker:2013cea} maps the
community program.  Each branch of this tree validates, at some point in its
development, against the two-dimensional calculation revisited here.

The audience for that calculation has changed.  Today the model is the
standard benchmark for Hamiltonian lattice methods, tensor networks, and
quantum simulation, and light-front quantization is itself being developed
as a basis for quantum
computation~\cite{Kreshchuk:2020dla,Kreshchuk:2020aiq,Li:2021kcs}.  Claims
of quantum advantage on 1+1d gauge theories are structurally weak, because
classical methods handle these theories so well~\cite{Lamm:2026vqz}.  This
paper sharpens that point with a ledger.  A direct qubit encoding of the
Fock space diagonalized here occupies $2N_c$ fermionic modes per momentum:
roughly 70 logical qubits to re-run the 1990 calculation, and 300--800 for
the resolutions and color groups in this paper.  The classical calculation
completes in seconds to minutes on one workstation.

The value of the old code to quantum simulation is therefore not as a
target but as a reference standard.  It supplies controlled,
error-budgeted continuum numbers against which equal-time hardware results
can be scored.  A reference is only as good as its validation, and DLCQ and
the equal-time formulations hardware uses are inequivalent quantizations.
Confirming the limits in which they agree is part of this paper's job.  The
one regime where they currently disagree is mapped in
Sec.~\ref{sec:modern}.

The model also has open questions of its own.  Light-front and equal-time
methods measure different chiral-limit behavior at finite $N$
(Sec.~\ref{sec:modern}).  Large-$N$ analyses predict a tetraquark bound
state near the two-meson threshold, and the prediction awaits a finite-$N$
test~\cite{Sazdjian:2025mqg,Lucha:2017gqq}.

This paper re-does the 1990 calculation with thirty-six years of hardware
and of error analysis.  The error analysis matters more.  We reproduce every figure and
the table of that paper, at the original parameters and again at
resolutions up to $2K \simeq 100$.  Two independent implementations of the
same Hamiltonian run against each other: the preserved Fortran of the
original program, and a validated open-source port.  The original
controlled its errors with one extrapolation series and a parenthesis
convention; we carry a four-component error budget with held-out validation
(Sec.~\ref{sec:budget}).  Its weak-coupling entries were limited by the
momentum grid in a way no resolution can repair.  We identify the
mechanism, quantify it without free parameters, and reach the chiral region
instead with van de Sande's improved
Hamiltonian~\cite{vandeSande:1996mc} (Sec.~\ref{sec:chiral}).

Two sections have no 1990 counterpart.  Section~\ref{sec:exotics} searches
for states dominated by four- and six-parton Fock sectors: tetraquark and
hexaquark candidates.  Section~\ref{sec:modern} compares against the modern
record: integrability at large $N$, tensor-network and quantum-hardware
spectra at $N = 2, 3$, and precision Schwinger-model benchmarks.

We keep the structure of Ref.~\cite{Hornbostel:1988fb} deliberately,
section for section where possible, so each figure can be laid beside its
1990 counterpart.  Sections~\ref{sec:lcq} and \ref{sec:sun} fix notation
and recall the Hamiltonian.  Section~\ref{sec:method} describes what
changed computationally.
Sections~\ref{sec:wavefunctions}--\ref{sec:massless} reproduce and extend
the 1990 results.  The remainder is new.
